% !TEX TS-program = lualatex

\documentclass[border = 3pt]{standalone}

%%%
% Des couleurs faciles d'emploi via le package \xcolor.
%%%
\usepackage[svgnames]{xcolor}

%%%
% La \bib \luadraw allie une facilité d’utilisation à un rendu
% particulièrement soigné.
%%%
\usepackage{luadraw}

%%%
% Externalisation de l'extraction des données pour permettre
% différents usages.
%%%
\directlua{dofile('../common/csv.lua')}

\begin{document}

\begin{luadraw*}{name = 3D-CSV-data-to-2D-repr}
------
-- Chargeons nos données (cf. path::''../common/csv.lua'').
------
local data = read_CSV(
  "../common/basic.csv",
-- On ignore la ¨1ere ligne qui contient les entêtes.
  true
)

------
-- ¨Def de la zone graphique.
------
local ld = luadraw

local graphview = ld.graph:new{
  window = {-7, 7, -5, 5},
  bbox   = false
}

------
-- Ajout des axes.
------
graphview:Daxes(
  {0, 1, 1},
  {
    arrows = "->",
    grid   = true
  }
)

------
-- Quelques réglages graphiques.
------
graphview:Linewidth(7)
graphview:Linecolor("Crimson")
graphview:Filloptions("full", "SandyBrown")

------
-- Traitement graphique des données.
------
local i = ld.cpx.I
local xyR

for _, xyR in ipairs(data) do
  graphview:Dcircle(xyR[1] + i*xyR[2], .5*xyR[3])
end

------
-- Montrons le résultat de notre oeuvre.
------
graphview:Show()
\end{luadraw*}

\end{document}
